% T7 — «Газета»: 2 колонки, засечный Times, узкие поля, диагональная подложка заголовка.
% Профиль: математика, планиметрия. 10 задач × 2 кол.
\documentclass[a4paper,10pt]{article}

\usepackage{../../../../Lessons/_templates/preamble_worksheet}
\usepackage{../_shared/worksheet_check}

% --- Узкие поля под «газетную» вёрстку ---
\geometry{a4paper, top=12mm, bottom=14mm, left=12mm, right=12mm}

% --- Палитра: глубокий тёмно-синий ---
\definecolor{wcprimary}{HTML}{1A365D}
\definecolor{wcprimaryLight}{HTML}{A0AEC0}
\definecolor{wcprimaryBG}{HTML}{EBF4FF}

% --- «Газетный» заголовок: подложка с поворотом + крупная засечная типографика ---
\renewcommand{\WorksheetTitle}[1]{%
  \par\noindent\begin{tikzpicture}
    % Диагональная цветная подложка
    \fill[wcprimary] (-1mm,-2mm) rectangle ([yshift=-14mm]\linewidth+1mm,2mm);
    \fill[wcprimary!75, rotate around={-1.2:(0,0)}]
      (-2mm,-15mm) rectangle (\linewidth+2mm, -1mm);
    % Лента-блик
    \fill[wcprimaryLight!50, rotate around={-1.2:(0,0)}]
      (-2mm,-15.5mm) rectangle (\linewidth+2mm, -14.5mm);
    % Текст заголовка
    \node[anchor=west, text=white, font=\fontsize{22pt}{26pt}\selectfont\bfseries]
      at (3mm, -7mm) {#1};
  \end{tikzpicture}\par\vspace{6mm}}

% --- TaskBox: газетная заметка — небольшой capital номер + текст ---
\renewcommand{\TaskBox}[2]{%
  \par\nopagebreak
  \noindent{\bfseries\color{wcprimary}\fontsize{11pt}{13pt}\selectfont\textsc{Задача №#1.}}\ %
  \small #2\par
  \vspace{0.5mm}{\color{wcprimaryLight}\rule{\linewidth}{0.3pt}}\par\vspace{2mm}}

\SetAnswerFieldSize{22mm}{7mm}

\begin{document}

\WorksheetTitle{Планиметрия. Треугольники}
\par{\small\itshape\color{wcprimary!80}\textsc{Вестник геометра} \quad·\quad ЕГЭ профиль, задание №1 \quad·\quad Класс 9 \quad·\quad T7 \quad·\quad ID: T7-planim-2026-05-26}\par\vspace{2mm}
\FirstPageHeader

\begin{multicols}{2}
\TaskBox{1}{В треугольнике $ABC$ угол $C$ равен $90^\circ$, $AC = 6$, $BC = 8$. Найдите гипотенузу.\par\hfill\textbf{Ответ:}~\AnswerField{1}}

\TaskBox{2}{Стороны треугольника равны $5$, $12$ и $13$. Найдите его площадь.\par\hfill\textbf{Ответ:}~\AnswerField{2}}

\TaskBox{3}{Один острый угол прямоугольного треугольника равен $25^\circ$. Найдите второй острый угол. Ответ в градусах.\par\hfill\textbf{Ответ:}~\AnswerField{3}}

\TaskBox{4}{Найдите площадь равнобедренного треугольника с основанием $12$ и боковой стороной $10$.\par\hfill\textbf{Ответ:}~\AnswerField{4}}

\TaskBox{5}{В треугольнике две стороны равны $7$ и $24$, а угол между ними прямой. Найдите медиану, проведённую к гипотенузе.\par\hfill\textbf{Ответ:}~\AnswerField{5}}

\TaskBox{6}{Сумма двух углов треугольника равна $135^\circ$. Найдите третий угол. Ответ в градусах.\par\hfill\textbf{Ответ:}~\AnswerField{6}}

\TaskBox{7}{Биссектриса прямого угла треугольника делит гипотенузу на отрезки $3$ и $4$. Найдите длину биссектрисы.\par\hfill\textbf{Ответ:}~\AnswerField{7}}

\TaskBox{8}{В равнобедренном треугольнике боковая сторона равна $13$, а основание $10$. Найдите высоту, проведённую к основанию.\par\hfill\textbf{Ответ:}~\AnswerField{8}}

\TaskBox{9}{Найдите радиус окружности, описанной около прямоугольного треугольника с катетами $9$ и $12$.\par\hfill\textbf{Ответ:}~\AnswerField{9}}

\TaskBox{10}{Внешний угол треугольника при одной из вершин равен $110^\circ$. Найдите сумму двух других углов, не смежных с этим внешним. Ответ в градусах.\par\hfill\textbf{Ответ:}~\AnswerField{10}}

\end{multicols}

\WorksheetQR{T7-planim/L1/v1}

\end{document}
